% !TEX program = xelatex
\documentclass[10pt,aspectratio=169]{ctexbeamer}

\usetheme[progressbar=frametitle,numbering=fraction]{metropolis}
\usepackage{appendixnumberbeamer}
\usepackage{booktabs}
\usepackage{multirow}
\usepackage{amsmath}
\usepackage{array}
\usepackage{tabularx}
\usepackage{graphicx}
\usepackage{tikz}
\usetikzlibrary{positioning,calc,fit,backgrounds,arrows.meta,shapes.geometric}
\usepackage{xcolor}
\usepackage{fontawesome5}

\setCJKmainfont[
    AutoFakeBold=true
]{FandolFang-Regular} % Beamer 默认使用 sf 字体族。

\setCJKsansfont[
    BoldFont=FandolSong-Bold,
    ItalicFont=FandolKai-Regular
]{FandolSong-Regular}

\setCJKmonofont[
    BoldFont=FandolHei-Bold
]{FandolHei-Regular}

\newCJKfontfamily\song[
    BoldFont=FandolSong-Bold
]{FandolSong-Regular}

\newCJKfontfamily\hei[
    BoldFont=FandolHei-Bold
]{FandolHei-Regular}

\newCJKfontfamily\kai[
    AutoFakeBold=true
]{FandolKai-Regular}

\newCJKfontfamily\fang[
    AutoFakeBold=true
]{FandolFang-Regular}


\definecolor{MainBlue}{HTML}{0F4C81}
\definecolor{AccentTeal}{HTML}{138A8A}
\definecolor{SoftGold}{HTML}{D89B2B}
\definecolor{LightBg}{HTML}{F5F7FA}
\definecolor{TextDark}{HTML}{1F2D3D}
\definecolor{SoftRed}{HTML}{C85C5C}

\setbeamercolor{normal text}{fg=TextDark,bg=white}
\setbeamercolor{frametitle}{fg=white,bg=MainBlue}
\setbeamercolor{title}{fg=white}
\setbeamercolor{subtitle}{fg=white}
\setbeamercolor{progress bar}{fg=SoftGold,bg=MainBlue!25}
\setbeamercolor{block title}{fg=white,bg=MainBlue}
\setbeamercolor{block body}{fg=TextDark,bg=LightBg}
\setbeamercolor{alertblock title}{fg=white,bg=SoftRed}
\setbeamercolor{alertblock body}{fg=TextDark,bg=SoftRed!8}
\setbeamercolor{exampleblock title}{fg=white,bg=AccentTeal}
\setbeamercolor{exampleblock body}{fg=TextDark,bg=AccentTeal!8}
\setbeamercolor{itemize item}{fg=MainBlue}
\setbeamercolor{itemize subitem}{fg=AccentTeal}
\setbeamercolor{section in toc}{fg=MainBlue}

\setbeamertemplate{footline}{
  \leavevmode
  \hbox{
    \begin{beamercolorbox}[wd=.82\paperwidth,ht=2.6ex,dp=1.2ex,leftskip=1em]{author in head/foot}
      \usebeamerfont{author in head/foot}智能交通课程汇报 \hfill STAEformer 模型研究与优化
    \end{beamercolorbox}
    \begin{beamercolorbox}[wd=.18\paperwidth,ht=2.6ex,dp=1.2ex,center]{date in head/foot}
      \insertframenumber/\inserttotalframenumber
    \end{beamercolorbox}
  }
}

\setbeamertemplate{headline}{}
\setbeamertemplate{navigation symbols}{}

\newcommand{\keymetric}[3]{
  \begin{tikzpicture}
    \node[
      rounded corners=8pt,
      fill=LightBg,
      draw=#1!70,
      line width=0.8pt,
      minimum width=3.1cm,
      minimum height=1.38cm,
      align=center
    ] {
      {\large\bfseries\color{#1} #2}\\[2pt]
      {\small #3}
    };
  \end{tikzpicture}
}

\newcommand{\featurepill}[2]{
  \begin{tikzpicture}[baseline=(x.base)]
    \node[
      rounded corners=10pt,
      fill=#1!12,
      draw=#1!70,
      line width=0.7pt,
      inner xsep=7pt,
      inner ysep=4pt
    ] (x) {\scriptsize\bfseries\color{#1} #2};
  \end{tikzpicture}
}

\newcommand{\sectionpagecontent}[1]{
  \begin{frame}[plain]
    \begin{tikzpicture}[remember picture,overlay]
      \fill[MainBlue] (current page.south west) rectangle (current page.north east);
      \fill[SoftGold] ([xshift=-1.2cm,yshift=1.2cm]current page.south east) circle (2.2cm);
      \fill[AccentTeal!80] ([xshift=1.0cm,yshift=-0.8cm]current page.north west) circle (1.7cm);
    \end{tikzpicture}
    \vspace*{2.2cm}
    {\color{white}\fontsize{24}{30}\selectfont\bfseries #1\par}
    \vspace{0.4cm}
    {\color{white!85}\large 面向交通流预测的模型分析、结构理解与实验总结}
  \end{frame}
}

\AtBeginSection{
  \sectionpagecontent{\insertsectionhead}
}

\title{交通流预测模型设计与优化}
\subtitle{基于 STAEformer 的调研、分析与实验汇报}
\author{智能交通课程作业}
\institute{交通预测方向专题报告}
\date{2026年3月25日}

\begin{document}

\begin{frame}[plain]
  \begin{tikzpicture}[remember picture,overlay]
    \fill[MainBlue] (current page.south west) rectangle (current page.north east);
    \fill[AccentTeal!85] ([xshift=-2.0cm,yshift=1.6cm]current page.south east) circle (3.0cm);
    \fill[SoftGold!92] ([xshift=1.3cm,yshift=-1.0cm]current page.north west) circle (2.1cm);
    \draw[white, line width=1.2pt, rounded corners=8pt]
      ([xshift=1.0cm,yshift=-1.1cm]current page.north west)
      rectangle
      ([xshift=-1.0cm,yshift=1.1cm]current page.south east);
  \end{tikzpicture}

  \vspace*{1.45cm}
  {\color{white}\fontsize{22}{26}\selectfont\bfseries 交通流预测模型设计与优化\par}
  \vspace{0.25cm}
  {\color{white!88}\Large 基于 STAEformer 的调研、分析与实验汇报\par}
  \vspace{0.65cm}

  \begin{columns}[T]
    \begin{column}{0.62\textwidth}
      \begin{block}{报告定位}
        \begin{itemize}
          \item 聚焦交通流预测任务中的时空依赖建模问题
          \item 从方法对比、模型结构、关键创新和实验结果展开
          \item 目标是给出“为什么选它、它怎么做、效果是否成立”的完整论证
        \end{itemize}
      \end{block}
    \end{column}
    \begin{column}{0.34\textwidth}
      \begin{exampleblock}{关键词}
        \centering
        Transformer\\[2pt]
        时空建模\\[2pt]
        自适应嵌入\\[2pt]
        多步预测
      \end{exampleblock}
    \end{column}
  \end{columns}

  \vspace{0.35cm}
  {\color{white!90}\normalsize 智能交通课程作业 \hfill 2026年3月25日}
\end{frame}

\begin{frame}{汇报目录}
  \tableofcontents[hideallsubsections]
\end{frame}

\section{研究背景与问题定义}

\begin{frame}{研究背景：为什么要做交通流预测}
  \begin{columns}[T]
    \begin{column}{0.60\textwidth}
      \begin{block}{任务背景}
        \begin{itemize}\setlength{\itemsep}{2pt}
          \item 交通流预测是智能交通系统的核心关键技术，直接支撑精细化管控、动态路径规划和出行时间估计
          \item 随着城市路网规模持续扩大，拥堵频发、效率下降等问题使高精度预测成为现实刚需
          \item 该任务本质上属于典型时空序列预测，需要同时处理历史依赖、节点关联和外部扰动
        \end{itemize}
      \end{block}

      \begin{alertblock}{应用价值}
        \begin{itemize}\setlength{\itemsep}{2pt}
          \item 为智能信号配时、诱导发布和拥堵预警提供前置决策依据
          \item 可预测流量、速度、占有率等关键交通状态量
          \item 兼具学术价值与工程部署意义，适合作为课程大作业研究对象
        \end{itemize}
      \end{alertblock}
    \end{column}
    \begin{column}{0.36\textwidth}
      \centering
      \keymetric{MainBlue}{多节点}{路网级联合预测}

      \vspace{0.18cm}
      \keymetric{AccentTeal}{多步长}{未来连续时刻输出}

      \vspace{0.18cm}
      \keymetric{SoftGold}{多因素}{时间与空间共同作用}
    \end{column}
  \end{columns}
\end{frame}

\begin{frame}{任务定义与评价目标}
  \begin{columns}[T]
    \begin{column}{0.53\textwidth}
      \begin{block}{预测任务定义}
        \small
        给定历史观测序列
        \[
          X \in \mathbb{R}^{B \times T_{in} \times N \times C},
        \]
        其中 $B$ 表示 batch，$T_{in}$ 表示历史时间步，$N$ 表示传感器节点数，$C$ 表示输入特征维度；
        模型输出未来多步预测
        \[
          \hat{Y} \in \mathbb{R}^{B \times T_{out} \times N \times 1}.
        \]
      \end{block}
      \begin{exampleblock}{核心挑战}
        \small 空间依赖复杂、时间模式多尺度、突发事件频繁，使得交通流预测必须同时具备精度、鲁棒性和可部署性。
      \end{exampleblock}
    \end{column}
    \begin{column}{0.43\textwidth}
      \begin{block}{评价目标}
        \begin{itemize}\setlength{\itemsep}{2pt}
          \item \textbf{MAE}：绝对误差，直观反映预测偏差
          \item \textbf{RMSE}：对大误差更敏感
          \item \textbf{MAPE}：便于衡量相对误差
        \end{itemize}
      \end{block}

      \begin{alertblock}{理想模型应满足}
        \begin{itemize}\setlength{\itemsep}{2pt}
          \item 准确率高
          \item 参数规模合理
          \item 训练稳定、易复现
          \item 对不同数据集有较好泛化性
        \end{itemize}
      \end{alertblock}
    \end{column}
  \end{columns}
\end{frame}

\begin{frame}{研究目标与章节安排}
  \begin{columns}[T]
    \begin{column}{0.58\textwidth}
      \begin{block}{研究目标}
        \begin{enumerate}\setlength{\itemsep}{3pt}
          \item 调研统计学习、RNN、GNN 与 Transformer 系列方法的发展脉络
          \item 以 STAEformer 为核心对象分析其嵌入机制、编码结构与预测流程
          \item 基于公开数据集完成模型复现，并用 MAE / RMSE / MAPE 验证结果
          \item 从长序列、异常事件、计算成本和外部因素建模等角度分析不足
          \item 在复现基础上提出改进模块，为后续研究提供可落地方向
        \end{enumerate}
      \end{block}
    \end{column}
    \begin{column}{0.38\textwidth}
      \begin{exampleblock}{报告主线}
        \featurepill{MainBlue}{调研}
        \featurepill{AccentTeal}{复现}
        \featurepill{SoftGold}{分析}

        \vspace{0.18cm}
        \featurepill{SoftRed}{改进}
        \featurepill{AccentTeal}{对比}
        \featurepill{MainBlue}{总结}
      \end{exampleblock}

      \begin{alertblock}{一句话概括}
        本次汇报采用“复现 - 分析 - 改进”的完整研究思路，而不是只介绍论文结论。
      \end{alertblock}
    \end{column}
  \end{columns}
\end{frame}

\section{方法调研与模型选型}

\begin{frame}{主流方法对比}
  \begin{table}
    \centering
    \scriptsize
    \renewcommand{\arraystretch}{1.18}
    \resizebox{0.98\linewidth}{!}{
    \begin{tabular}{>{\raggedright\arraybackslash}p{1.6cm} >{\raggedright\arraybackslash}p{2.0cm} >{\centering\arraybackslash}p{1.7cm} >{\centering\arraybackslash}p{1.7cm} >{\raggedright\arraybackslash}p{2.8cm}}
      \toprule
      \textbf{方法类别} & \textbf{代表模型} & \textbf{时间建模} & \textbf{空间建模} & \textbf{主要局限} \\
      \midrule
      CNN 类 & TCN / STGCN & 局部卷积 & 依赖预定义图 & 感受野有限，长依赖建模不充分 \\
      RNN 类 & LSTM / GRU & 递归记忆 & 弱或间接建模 & 串行计算，长序列训练效率低 \\
      GNN 类 & DCRNN / Graph WaveNet & 较强 & 显式图传播 & 需要图结构或复杂传播机制 \\
      Transformer 类 & GMAN / ASTTN / STAEformer & 全局注意力 & 可学习或分离建模 & 需控制计算复杂度 \\
      \bottomrule
    \end{tabular}
    }
  \end{table}

  \begin{columns}[T]
    \begin{column}{0.52\textwidth}
      \begin{block}{选择 Transformer 路线的原因}
        \begin{itemize}
          \item 更擅长捕获长程依赖，适合多步交通预测
          \item 并行计算友好，训练效率优于纯 RNN 结构
          \item 模型表达能力强，便于统一时间和空间信息
        \end{itemize}
      \end{block}
    \end{column}
    \begin{column}{0.44\textwidth}
      \begin{exampleblock}{选型结论}
        \featurepill{MainBlue}{长程依赖}
        \featurepill{AccentTeal}{并行训练}
        \featurepill{SoftGold}{结构清晰}

        \vspace{0.15cm}
        \small STAEformer 同时兼顾效果、结构简洁性和复现可行性，因此被选为本次分析对象。
      \end{exampleblock}
    \end{column}
  \end{columns}
\end{frame}

\begin{frame}{为什么选择 STAEformer}
  \begin{columns}[T]
    \begin{column}{0.58\textwidth}
      \begin{block}{模型优势}
        \begin{enumerate}
          \item 不强依赖人工构造图结构，降低先验设计负担
          \item 将时空注意力拆分，缓解直接联合注意力带来的高复杂度
          \item 自适应嵌入可以直接学习“节点-时间”级别的表示
          \item 结构比较清晰，便于解释模型每个模块的功能
        \end{enumerate}
      \end{block}
    \end{column}
    \begin{column}{0.38\textwidth}
      \begin{alertblock}{适合作业汇报的原因}
        \begin{itemize}
          \item 既有理论亮点
          \item 又能清楚展示代码实现
          \item 实验结果易于量化呈现
        \end{itemize}
      \end{alertblock}

      \begin{exampleblock}{一句话概括}
        STAEformer 的核心思路是：用可学习嵌入替代固定图先验，用时空分离注意力提高效率。
      \end{exampleblock}
    \end{column}
  \end{columns}
\end{frame}

\begin{frame}{选题模型概述}
  \begin{columns}[T]
    \begin{column}{0.57\textwidth}
      \begin{block}{STAEformer 的核心思想}
        STAEformer 全称为 \textbf{Spatio-Temporal Adaptive Embedding Transformer}。其设计重点不是堆叠复杂模块，而是先构造高质量时空表示，再交给较简洁的 Transformer 骨干完成建模。
      \end{block}

      \begin{exampleblock}{与常见复杂模型的区别}
        \begin{itemize}\setlength{\itemsep}{2pt}
          \item 不强调图卷积、门控模块和多分支结构的反复叠加
          \item 更强调“高质量时空表示 + 简洁骨干”的范式
          \item 通过统一嵌入将空间依赖与时间动态隐式融合到同一特征空间
        \end{itemize}
      \end{exampleblock}
    \end{column}
    \begin{column}{0.39\textwidth}
      \begin{alertblock}{模型定位}
        \begin{itemize}\setlength{\itemsep}{2pt}
          \item 预测精度高
          \item 结构规整清晰
          \item 实现难度可控
          \item 复现与部署友好
        \end{itemize}
      \end{alertblock}

      \begin{block}{结论}
        借助高质量前期嵌入，Vanilla Transformer 也能有效建模复杂交通时空依赖，并实现较好的性能与工程平衡。
      \end{block}
    \end{column}
  \end{columns}
\end{frame}

\section{模型结构与关键机制}

\begin{frame}{STAEformer 整体框架}
  \centering
  \resizebox{0.96\linewidth}{!}{
  \begin{tikzpicture}[
    node distance=0.55cm and 0.7cm,
    module/.style={rounded corners=6pt, minimum width=2.0cm, minimum height=1.0cm, align=center, draw, line width=0.8pt},
    arrow/.style={-{Latex[length=3mm]}, thick, MainBlue}
  ]
    \node[module, fill=AccentTeal!18, draw=AccentTeal] (input) {输入序列\\$X(B,T,N,C)$};
    \node[module, fill=MainBlue!10, draw=MainBlue, right=of input] (proj) {线性映射\\Input Projection};
    \node[module, fill=SoftGold!20, draw=SoftGold, right=of proj] (embed) {时空自适应\\嵌入融合};
    \node[module, fill=AccentTeal!18, draw=AccentTeal, right=of embed] (temp) {Temporal\\Attention};
    \node[module, fill=MainBlue!12, draw=MainBlue, right=of temp] (spatial) {Spatial\\Attention};
    \node[module, fill=SoftGold!20, draw=SoftGold, right=of spatial] (ffn) {FFN\\前馈网络};
    \node[module, fill=SoftRed!15, draw=SoftRed, right=of ffn] (out) {Mixed\\Projection};
    \node[module, fill=green!18!white, draw=green!50!black, right=of out] (pred) {预测输出\\$\hat{Y}$};

    \draw[arrow] (input) -- (proj);
    \draw[arrow] (proj) -- (embed);
    \draw[arrow] (embed) -- (temp);
    \draw[arrow] (temp) -- (spatial);
    \draw[arrow] (spatial) -- (ffn);
    \draw[arrow] (ffn) -- (out);
    \draw[arrow] (out) -- (pred);
  \end{tikzpicture}
  }

  \vspace{0.3cm}
  \begin{columns}[T]
    \begin{column}{0.48\textwidth}
      \begin{block}{输入输出维度}
        历史输入通常为 $(B,12,N,3)$，未来输出为 $(B,12,N,1)$，适用于多节点、多步长预测场景。
      \end{block}
    \end{column}
    \begin{column}{0.48\textwidth}
      \begin{exampleblock}{结构特点}
        模型主干保持简洁，把复杂性集中在“嵌入设计”和“注意力分解”上，解释性较好。
      \end{exampleblock}
    \end{column}
  \end{columns}
\end{frame}

\begin{frame}{关键创新一：时空自适应嵌入}
  \begin{columns}[T]
    \begin{column}{0.52\textwidth}
      \begin{block}{嵌入组成}
        \begin{table}
          \centering
          \scriptsize
          \renewcommand{\arraystretch}{1.14}
          \resizebox{0.98\linewidth}{!}{
          \begin{tabular}{lcc}
            \toprule
            \textbf{嵌入类型} & \textbf{维度} & \textbf{作用} \\
            \midrule
            Input Embedding & 24 & 原始特征映射 \\
            ToD Embedding & 24 & 时间片信息编码 \\
            DoW Embedding & 24 & 星期模式编码 \\
            Adaptive Embedding & 80 & 节点-时间联合表示 \\
            \midrule
            总维度 & 152 & 统一输入表示 \\
            \bottomrule
          \end{tabular}
          }
        \end{table}
      \end{block}
    \end{column}
    \begin{column}{0.44\textwidth}
      \begin{alertblock}{设计意义}
        \begin{itemize}
          \item 不再完全依赖静态邻接矩阵
          \item 让模型自行学习不同节点在不同时间片的潜在关系
          \item 提升对复杂交通模式的表达能力
        \end{itemize}
      \end{alertblock}

      \begin{exampleblock}{直观理解}
        固定图结构刻画的是“谁和谁可能相关”，自适应嵌入更进一步学习“在什么时间、相关到什么程度”。
      \end{exampleblock}
    \end{column}
  \end{columns}
\end{frame}

\begin{frame}{自适应嵌入与图神经网络对比}
  \begin{table}
    \centering
    \scriptsize
    \renewcommand{\arraystretch}{1.2}
    \resizebox{0.96\linewidth}{!}{
    \begin{tabular}{>{\raggedright\arraybackslash}p{1.9cm} >{\centering\arraybackslash}p{2.7cm} >{\centering\arraybackslash}p{2.8cm} >{\centering\arraybackslash}p{2.0cm}}
      \toprule
      \textbf{方面} & \textbf{图神经网络} & \textbf{STAEformer} & \textbf{说明} \\
      \midrule
      图结构依赖 & 需要预定义邻接矩阵 & 不需要 & 减轻图先验设计负担 \\
      空间建模方式 & 基于图结构稀疏连接 & 全连接自适应关系 & 更偏数据驱动 \\
      灵活性 & 受固定图结构约束 & 可学习节点-时间表示 & 对复杂关联更灵活 \\
      复杂度 & $\mathcal{O}(E \times D)$ & $\mathcal{O}(N^2 \times D)$ & 大规模路网下需关注开销 \\
      \bottomrule
    \end{tabular}
    }
  \end{table}

  \vspace{0.15cm}
  \begin{columns}[T]
    \begin{column}{0.48\textwidth}
      \begin{exampleblock}{优势}
        不需要手工给定固定图结构，能够从数据中直接学习潜在时空关系。
      \end{exampleblock}
    \end{column}
    \begin{column}{0.48\textwidth}
      \begin{alertblock}{代价}
        当节点数增大时，空间注意力复杂度仍然较高，这也是后续改进的重要出发点。
      \end{alertblock}
    \end{column}
  \end{columns}
\end{frame}

\begin{frame}{关键创新二：时空分离注意力}
  \begin{columns}[T]
    \begin{column}{0.54\textwidth}
      \begin{block}{核心思路}
        \small
        如果直接对 $(T \times N)$ 个时空 token 做联合注意力，计算和显存开销会迅速增大。STAEformer 将其拆分为：
        \begin{itemize}
          \item 先沿时间维度做 Temporal Attention
          \item 再沿空间维度做 Spatial Attention
        \end{itemize}
        这样在保留建模能力的同时，能显著降低复杂度。
      \end{block}

      \begin{exampleblock}{复杂度对比}
        联合时空注意力约为 $\mathcal{O}((TN)^2)$，分离后可近似写作 $\mathcal{O}(T^2 + N^2)$，更适合实际训练。
      \end{exampleblock}
    \end{column}
    \begin{column}{0.42\textwidth}
      \centering
      \resizebox{0.9\linewidth}{!}{
      \begin{tikzpicture}[
        box/.style={rounded corners=5pt, draw, minimum width=2.9cm, minimum height=0.9cm, align=center},
        line/.style={-{Latex[length=2.5mm]}, thick}
      ]
        \node[box, fill=MainBlue!10, draw=MainBlue] (x) {输入表示 $H$};
        \node[box, fill=AccentTeal!15, draw=AccentTeal, below=0.45cm of x] (t) {Temporal Attention};
        \node[box, fill=SoftGold!18, draw=SoftGold, below=0.45cm of t] (s) {Spatial Attention};
        \node[box, fill=SoftRed!12, draw=SoftRed, below=0.45cm of s] (y) {输出表示};
        \draw[line] (x) -- (t);
        \draw[line] (t) -- (s);
        \draw[line] (s) -- (y);
      \end{tikzpicture}
      }
    \end{column}
  \end{columns}
\end{frame}

\begin{frame}{STAE Layer 内部结构}
  \begin{center}
    \resizebox{0.96\linewidth}{!}{
    \begin{tikzpicture}[
      node distance=0.38cm and 0.55cm,
      unit/.style={rounded corners=5pt, draw, minimum width=1.8cm, minimum height=0.82cm, align=center},
      arrow/.style={-{Latex[length=2.8mm]}, thick}
    ]
      \node[unit, fill=gray!12] (x) {输入 $x$};
      \node[unit, fill=MainBlue!12, draw=MainBlue, right=of x] (ln1) {LayerNorm};
      \node[unit, fill=AccentTeal!16, draw=AccentTeal, right=of ln1] (ta) {Temporal\\Attention};
      \node[unit, fill=MainBlue!12, draw=MainBlue, right=of ta] (ln2) {LayerNorm};
      \node[unit, fill=SoftGold!18, draw=SoftGold, right=of ln2] (sa) {Spatial\\Attention};
      \node[unit, fill=MainBlue!12, draw=MainBlue, right=of sa] (ln3) {LayerNorm};
      \node[unit, fill=SoftRed!12, draw=SoftRed, right=of ln3] (ffn) {Feed\\Forward};
      \node[unit, fill=gray!12, right=of ffn] (y) {输出 $y$};

      \draw[arrow] (x) -- (ln1);
      \draw[arrow] (ln1) -- (ta);
      \draw[arrow] (ta) -- (ln2);
      \draw[arrow] (ln2) -- (sa);
      \draw[arrow] (sa) -- (ln3);
      \draw[arrow] (ln3) -- (ffn);
      \draw[arrow] (ffn) -- (y);

      \draw[dashed, thick, MainBlue] ($(x.north)+(0,0.38)$) to[bend left=18] ($(ta.north)+(0,0.38)$);
      \draw[dashed, thick, SoftGold!80!black] ($(ta.south)+(0,-0.36)$) to[bend right=18] ($(sa.south)+(0,-0.36)$);
      \draw[dashed, thick, SoftRed!85!black] ($(sa.north)+(0,0.38)$) to[bend left=18] ($(y.north)+(0,0.38)$);
    \end{tikzpicture}
    }
  \end{center}

  \vspace{0.35cm}
  \begin{block}{模块理解}
    一个 STAE Layer 可以看成“归一化 + 时间建模 + 空间建模 + 非线性映射”的串联模块，多层堆叠后逐步提升时空表征能力。
  \end{block}
\end{frame}

\begin{frame}{输出层：Mixed Projection 的作用}
  \begin{columns}[T]
    \begin{column}{0.54\textwidth}
      \begin{block}{为什么不用逐步自回归输出}
        传统逐步预测往往误差会累积，而 Mixed Projection 直接将隐藏表示映射为未来多个时间步，能够更稳定地完成多步预测。
      \end{block}
      \begin{exampleblock}{优势}
        \begin{itemize}
          \item 端到端输出未来序列
          \item 减少滚动预测带来的误差传播
          \item 结构简单，和主干模块衔接自然
        \end{itemize}
      \end{exampleblock}
    \end{column}
    \begin{column}{0.42\textwidth}
      \centering
      \resizebox{0.9\linewidth}{!}{
      \begin{tikzpicture}[
        box/.style={rounded corners=5pt, draw, minimum width=3.1cm, minimum height=0.82cm, align=center},
        arrow/.style={-{Latex[length=2.6mm]}, thick}
      ]
        \node[box, fill=MainBlue!10, draw=MainBlue] (a) {隐藏表示 $H$};
        \node[box, fill=gray!10, below=0.42cm of a] (b) {重排 / reshape};
        \node[box, fill=SoftGold!18, draw=SoftGold, below=0.42cm of b] (c) {线性映射到 $T_{out}$};
        \node[box, fill=AccentTeal!15, draw=AccentTeal, below=0.42cm of c] (d) {预测输出 $\hat{Y}$};
        \draw[arrow] (a) -- (b);
        \draw[arrow] (b) -- (c);
        \draw[arrow] (c) -- (d);
      \end{tikzpicture}
      }
    \end{column}
  \end{columns}
\end{frame}

\section{实验设置与结果分析}

\begin{frame}{实验设计与评价指标}
  \begin{columns}[T]
    \begin{column}{0.56\textwidth}
      \begin{block}{实验设计}
        \begin{itemize}\setlength{\itemsep}{2pt}
          \item 使用公开交通数据集开展复现实验，核心验证模型短时多步预测能力
          \item 通过定量指标、定性曲线和工程开销三条线共同评价模型
          \item 在复现 STAEformer 基线后，再引入改进模块做横向对比
        \end{itemize}
      \end{block}

      \begin{exampleblock}{结果分析维度}
        \begin{itemize}\setlength{\itemsep}{2pt}
          \item 定量分析：MAE / RMSE / MAPE 的整体对比
          \item 定性分析：高峰期、平峰期与突变点的拟合表现
          \item 工程分析：参数量、显存占用、训练耗时和推理速度
        \end{itemize}
      \end{exampleblock}
    \end{column}
    \begin{column}{0.40\textwidth}
      \begin{alertblock}{指标说明}
        \begin{itemize}\setlength{\itemsep}{2pt}
          \item \textbf{MAE}：整体偏差越小越好
          \item \textbf{RMSE}：更关注极端误差
          \item \textbf{MAPE}：便于跨模型跨数据集比较
        \end{itemize}
      \end{alertblock}

      \begin{block}{预期验证结论}
        STAEformer 应在短时预测精度、收敛效率和结构简洁性上体现较强综合优势。
      \end{block}
    \end{column}
  \end{columns}
\end{frame}

\begin{frame}{实验配置概览}
  \begin{columns}[T]
    \begin{column}{0.50\textwidth}
      \begin{block}{主要超参数}
        \begin{table}
          \centering
          \scriptsize
          \renewcommand{\arraystretch}{1.12}
          \begin{tabular}{lc}
            \toprule
            \textbf{参数} & \textbf{取值} \\
            \midrule
            num\_nodes & 207 \\
            in\_steps / out\_steps & 12 / 12 \\
            input\_dim / output\_dim & 3 / 1 \\
            model\_dim & 152 \\
            num\_heads & 4 \\
            num\_layers & 3 \\
            feed\_forward\_dim & 256 \\
            adaptive\_embedding\_dim & 80 \\
            \bottomrule
          \end{tabular}
        \end{table}
      \end{block}
    \end{column}
    \begin{column}{0.46\textwidth}
      \begin{alertblock}{配置特点}
        \begin{itemize}
          \item 模型容量适中
          \item 不依赖超大显存即可训练
          \item 结构设计与论文实现保持一致
        \end{itemize}
      \end{alertblock}

      \begin{exampleblock}{说明}
        \featurepill{MainBlue}{207 节点}
        \featurepill{AccentTeal}{3 层编码}
        \featurepill{SoftGold}{4 头注意力}

        \vspace{0.18cm}
        这类配置适合课程作业场景，既有研究价值，也具有较强可复现性。
      \end{exampleblock}
    \end{column}
  \end{columns}
\end{frame}

\begin{frame}{模型复现：当前已完成内容}
  \begin{columns}[T]
    \begin{column}{0.56\textwidth}
      \begin{block}{基线复现进展}
        \begin{itemize}\setlength{\itemsep}{2pt}
          \item 已完成 STAEformer 基线训练脚本梳理与参数核对
          \item 本地改进实验框架已搭建，包含 Baseline、EFF、DGL、LA、MST 五类模型
          \item 当前仓库中已有 \textbf{PEMS08} 基线训练日志，可确认数据划分、模型参数量和训练入口已跑通
        \end{itemize}
      \end{block}

      \begin{exampleblock}{PEMS08 已确认信息}
        Train: 10700，Val: 3567，Test: 3566；输入窗口 12，输出窗口 12；总参数量 \textbf{1,223,460}。
      \end{exampleblock}
    \end{column}
    \begin{column}{0.40\textwidth}
      \begin{alertblock}{已完成证据}
        \begin{itemize}\setlength{\itemsep}{2pt}
          \item \texttt{baseline\_PEMS08\_train.log}
          \item \texttt{train\_baseline.py}
          \item \texttt{train.py}
          \item 各改进模型 \texttt{model.py}
        \end{itemize}
      \end{alertblock}

      \begin{block}{说明}
        目前仓库里尚未生成各模型完整 \texttt{json} 结果文件，因此本页只展示已确认复现数据，不虚构最终误差指标。
      \end{block}
    \end{column}
  \end{columns}
\end{frame}

\begin{frame}{复现实验数据与工程信息}
  \begin{table}
    \centering
    \scriptsize
    \renewcommand{\arraystretch}{1.18}
    \resizebox{0.96\linewidth}{!}{
    \begin{tabular}{>{\centering\arraybackslash}p{1.6cm} >{\centering\arraybackslash}p{1.7cm} >{\centering\arraybackslash}p{1.5cm} >{\centering\arraybackslash}p{1.6cm} >{\centering\arraybackslash}p{1.8cm} >{\centering\arraybackslash}p{1.8cm}}
      \toprule
      \textbf{数据集} & \textbf{节点数} & \textbf{采样间隔} & \textbf{输入/输出步长} & \textbf{本地已确认信息} & \textbf{备注} \\
      \midrule
      PEMS08 & 170 & 5 min & 12 / 12 & Train 10700 / Val 3567 / Test 3566 & 已有训练日志 \\
      METR-LA & 207 & 5 min & 12 / 12 & 参数配置已整理 & 待补复现结果 \\
      PEMS-BAY & 325 & 5 min & 12 / 12 & 参数配置已整理 & 待补复现结果 \\
      \bottomrule
    \end{tabular}
    }
  \end{table}

  \vspace{0.15cm}
  \begin{columns}[T]
    \begin{column}{0.48\textwidth}
      \begin{exampleblock}{工程侧观察}
        STAEformer 基线在 PEMS08 上参数量约 1.22M，属于较易复现、便于后续做模块化改进的骨干模型。
      \end{exampleblock}
    \end{column}
    \begin{column}{0.48\textwidth}
      \begin{alertblock}{汇报建议}
        口头说明时可强调：本次作业不是停留在阅读论文，而是已经把复现入口、配置和改进框架都搭了起来。
      \end{alertblock}
    \end{column}
  \end{columns}
\end{frame}

\begin{frame}{复现结果展示页设计}
  \begin{columns}[T]
    \begin{column}{0.52\textwidth}
      \begin{block}{基线模型结果表}
        \begin{table}
          \centering
          \scriptsize
          \renewcommand{\arraystretch}{1.15}
          \begin{tabular}{lccc}
            \toprule
            \textbf{模型} & \textbf{MAE} & \textbf{RMSE} & \textbf{MAPE} \\
            \midrule
            STAEformer Baseline & 待补 & 待补 & 待补 \\
            \bottomrule
          \end{tabular}
        \end{table}
      \end{block}

      \begin{exampleblock}{当前状态}
        由于结果 \texttt{json} 尚未生成，误差指标位先留作实验完成后的填充位置；表格结构已经留好，可直接替换。
      \end{exampleblock}
    \end{column}
    \begin{column}{0.44\textwidth}
      \begin{block}{本页建议讲法}
        \begin{itemize}\setlength{\itemsep}{2pt}
          \item 先说明复现环境、数据集和训练配置
          \item 再展示基线误差指标与训练开销
          \item 最后引出“为什么还需要做改进”
        \end{itemize}
      \end{block}

      \begin{alertblock}{过渡作用}
        这一页的重点是把“复现成功”与“继续优化”顺畅衔接起来。
      \end{alertblock}
    \end{column}
  \end{columns}
\end{frame}

\begin{frame}{结果解读：模型表现为什么更好}
  \begin{columns}[T]
    \begin{column}{0.48\textwidth}
      \begin{block}{性能优势来源}
        \begin{itemize}
          \item 自适应嵌入提高了对复杂空间关系的表达能力
          \item 分离注意力兼顾了时间依赖和空间依赖
          \item 多步直接预测减少了误差累积
          \item 结构清晰，训练过程中更稳定
        \end{itemize}
      \end{block}
    \end{column}
    \begin{column}{0.48\textwidth}
      \begin{block}{实验汇报可强调的观察}
        \begin{itemize}
          \item 短时预测和中时预测都保持较好精度
          \item 在多节点路网场景中泛化性更强
          \item 相较传统图模型，对图先验依赖更弱
          \item 在效果和实现复杂度之间取得更好平衡
        \end{itemize}
      \end{block}
    \end{column}
  \end{columns}

  \vspace{0.1cm}
  \begin{alertblock}{答辩表达建议}
    如果老师更关注“你做了什么”，这一页可以口头补充：代码复现、参数理解、模块拆解、误差指标对比和模型优缺点分析。
  \end{alertblock}
\end{frame}

\begin{frame}{模型不足与问题定位}
  \begin{columns}[T]
    \begin{column}{0.48\textwidth}
      \begin{alertblock}{当前局限}
        \begin{itemize}
          \item 长序列预测场景下误差仍可能逐步积累
          \item 极端突发事件下仍可能预测失稳
          \item 自适应嵌入可解释性仍然有限
          \item 大规模路网下空间注意力计算成本较高
        \end{itemize}
      \end{alertblock}
    \end{column}
    \begin{column}{0.48\textwidth}
      \begin{exampleblock}{后续优化方向}
        \begin{itemize}
          \item 融合天气、事件、节假日等外部因素
          \item 引入动态图或在线更新机制增强空间依赖刻画
          \item 使用轻量化注意力降低大图计算开销
          \item 采用多尺度时间建模提升长周期建模能力
        \end{itemize}
      \end{exampleblock}
    \end{column}
  \end{columns}

  \vspace{0.25cm}
  \begin{block}{总结判断}
    STAEformer 不是简单地“把 Transformer 用在交通预测上”，而是围绕交通数据的时空特性做了结构化改造，因此更适合作为课程汇报中的重点分析对象。
  \end{block}
\end{frame}

\section{改进方案设计}

\begin{frame}{改进方案总览}
  \begin{columns}[T]
    \begin{column}{0.58\textwidth}
      \begin{block}{四类改进模块}
        \begin{enumerate}\setlength{\itemsep}{3pt}
          \item \textbf{EFF}：外部特征融合，引入天气、节假日、道路事件等信息
          \item \textbf{DGL}：动态图结构学习，增强节点间动态空间依赖建模
          \item \textbf{LA}：轻量化注意力，降低大规模路网下的计算开销
          \item \textbf{MST}：多尺度时间建模，联合建模短期、日周期和周周期
        \end{enumerate}
      \end{block}

      \begin{exampleblock}{总体原则}
        在不破坏 STAEformer “高质量时空表示 + 简洁 Transformer 骨干” 设计范式的前提下，采用模块化方式嵌入改进组件。
      \end{exampleblock}
    \end{column}
    \begin{column}{0.38\textwidth}
      \begin{alertblock}{改进目标}
        \begin{itemize}\setlength{\itemsep}{2pt}
          \item 更强异常场景适应能力
          \item 更灵活的空间依赖建模
          \item 更低推理成本
          \item 更稳的长序列预测表现
        \end{itemize}
      \end{alertblock}
    \end{column}
  \end{columns}
\end{frame}

\begin{frame}{四类改进模块的实现思路}
  \begin{table}
    \centering
    \scriptsize
    \renewcommand{\arraystretch}{1.18}
    \resizebox{0.98\linewidth}{!}{
    \begin{tabular}{>{\centering\arraybackslash}p{1.0cm} >{\centering\arraybackslash}p{2.3cm} >{\centering\arraybackslash}p{4.2cm} >{\centering\arraybackslash}p{2.6cm}}
      \toprule
      \textbf{简称} & \textbf{改进方向} & \textbf{实现要点} & \textbf{主要收益} \\
      \midrule
      EFF & 外部特征融合 & 将天气、节假日、事件等特征编码后与原始交通状态共同输入嵌入层 & 强化异常场景鲁棒性 \\
      DGL & 动态图学习 & 基于实时交通状态动态生成邻接关系或图结构表示 & 增强空间依赖灵活性 \\
      LA & 轻量化注意力 & 用稀疏、线性或分块注意力替代稠密注意力 & 降低复杂度与推理成本 \\
      MST & 多尺度时间建模 & 分别编码短期、日周期、周周期信息后再做特征融合 & 提升长时预测稳定性 \\
      \bottomrule
    \end{tabular}
    }
  \end{table}

  \vspace{0.15cm}
  \begin{alertblock}{实现方式}
    改进模块均已在 `improvements/` 目录中建立独立实现入口，便于和基线模型做公平对比。
  \end{alertblock}
\end{frame}

\begin{frame}{改进后的整体实现思路}
  \centering
  \resizebox{0.95\linewidth}{!}{
  \begin{tikzpicture}[
    node distance=0.4cm and 0.55cm,
    box/.style={rounded corners=6pt, minimum width=2.05cm, minimum height=0.9cm, align=center, draw, line width=0.8pt},
    arrow/.style={-{Latex[length=2.8mm]}, thick, MainBlue}
  ]
    \node[box, fill=SoftGold!18, draw=SoftGold] (eff) {外部特征融合\\EFF};
    \node[box, fill=AccentTeal!14, draw=AccentTeal, right=of eff] (mst) {多尺度时间\\编码 MST};
    \node[box, fill=MainBlue!10, draw=MainBlue, right=of mst] (embed) {时空自适应嵌入\\+ 动态图学习 DGL};
    \node[box, fill=SoftRed!12, draw=SoftRed, right=of embed] (attn) {轻量化注意力\\Transformer};
    \node[box, fill=green!18!white, draw=green!50!black, right=of attn] (out) {预测输出};
    \draw[arrow] (eff) -- (mst);
    \draw[arrow] (mst) -- (embed);
    \draw[arrow] (embed) -- (attn);
    \draw[arrow] (attn) -- (out);
  \end{tikzpicture}
  }

  \vspace{0.25cm}
  \begin{block}{设计原则}
    输入端新增外部特征融合与多尺度时间编码，嵌入层增强动态空间关系学习，编码层用轻量化注意力替代稠密注意力，输出层保持不变。
  \end{block}
\end{frame}

\section{改进结果与对比}

\begin{frame}{改进实验结果表设计}
  \begin{columns}[T]
    \begin{column}{0.58\textwidth}
      \begin{block}{模型对比表}
        \begin{table}
          \centering
          \scriptsize
          \renewcommand{\arraystretch}{1.15}
          \begin{tabular}{lccccc}
            \toprule
            \textbf{模型} & \textbf{MAE} & \textbf{RMSE} & \textbf{MAPE} & \textbf{参数量} & \textbf{推理} \\
            \midrule
            Baseline & 待补 & 待补 & 待补 & 1.22M & 待补 \\
            EFF & 待补 & 待补 & 待补 & $\sim$1.3M & 待补 \\
            DGL & 待补 & 待补 & 待补 & $\sim$1.4M & 待补 \\
            LA & 待补 & 待补 & 待补 & $\sim$1.2M & 待补 \\
            MST & 待补 & 待补 & 待补 & $\sim$1.3M & 待补 \\
            \bottomrule
          \end{tabular}
        \end{table}
      \end{block}
    \end{column}
    \begin{column}{0.38\textwidth}
      \begin{alertblock}{本页用途}
        \begin{itemize}\setlength{\itemsep}{2pt}
          \item 用于展示最终实验结果
          \item 当前参数量可先填入工程估计值
        \item 指标结果待 \texttt{results/*.json} 生成后补齐
        \end{itemize}
      \end{alertblock}

      \begin{exampleblock}{结论模板}
        可以按“谁精度最好、谁速度最快、谁最适合部署”三个角度总结改进效果。
      \end{exampleblock}
    \end{column}
  \end{columns}
\end{frame}

\begin{frame}{改进结果的分析口径}
  \begin{columns}[T]
    \begin{column}{0.48\textwidth}
      \begin{block}{若结果显示 EFF 最优}
        说明外部因素确实对异常场景和高波动时段具有显著补充作用。
      \end{block}

      \begin{block}{若结果显示 DGL 最优}
        说明静态嵌入之外，动态空间结构学习对复杂路网建模仍有增益。
      \end{block}
    \end{column}
    \begin{column}{0.48\textwidth}
      \begin{block}{若结果显示 LA 最快}
        说明轻量化注意力在大规模路网部署中更具工程价值。
      \end{block}

      \begin{block}{若结果显示 MST 更稳}
        说明多尺度时间编码有助于提升长时依赖建模和长序列预测稳定性。
      \end{block}
    \end{column}
  \end{columns}

  \vspace{0.15cm}
  \begin{alertblock}{答辩建议}
    如果最终只跑出部分实验，也可以明确说明：已完成基线复现，改进模型框架已实现，当前重点展示的是设计逻辑、对比框架与可补充的结果页。
  \end{alertblock}
\end{frame}

\section{总结}

\begin{frame}[plain]
  \begin{tikzpicture}[remember picture,overlay]
    \fill[MainBlue] (current page.south west) rectangle (current page.north east);
    \fill[AccentTeal!85] ([xshift=-1.5cm,yshift=1.1cm]current page.south east) circle (2.3cm);
    \fill[SoftGold!95] ([xshift=1.3cm,yshift=-0.9cm]current page.north west) circle (1.8cm);
  \end{tikzpicture}

  \vspace*{2.1cm}
  \begin{center}
    {\color{white}\fontsize{24}{30}\selectfont\bfseries 汇报总结\par}
    \vspace{0.55cm}
    {\color{white!92}\Large STAEformer 通过时空自适应嵌入与分离注意力机制，\\实现了对交通流时空依赖的高效建模。\par}
    \vspace{0.9cm}
    {\color{white!85}\large 关键词：结构清晰、效果优秀、适合复现、便于讲解}
  \end{center}
\end{frame}

\begin{frame}[standout]
  谢谢观看
\end{frame}

\end{document}
